\documentclass[12pt]{article}
\usepackage{amsmath, amssymb, amsthm, geometry, setspace}
\geometry{margin=1in}
\setstretch{1.25}
\setlength{\parindent}{0pt}

\title{Quiz 6\\ \vspace{5pt}
\large ECON 320 -- Introduction to Mathematical Economics}
\author{Instructor: Muhebullah Karimzada}
\date{December 01, 2025}

\begin{document}
\maketitle

\section*{Part A — Multiple Choice (1 mark each)}

\begin{enumerate}

\item The Envelope Theorem tells us that when evaluating the effect of a parameter on the value function:
\begin{enumerate}
\item We must take into account the indirect effect through optimal choices.
\item We ignore all partial derivatives involving the parameter.
\item Only the direct partial effect matters; indirect effects cancel by FOCs.
\item The derivative of the parameter is always zero.
\end{enumerate}

\item In constrained optimization, the Lagrange multiplier represents:
\begin{enumerate}
\item The marginal cost of production.
\item The shadow value of relaxing the constraint by one unit.
\item The total value of the objective function.
\item A second-order condition checker.
\end{enumerate}

\item The indirect utility function is the \textbf{dual} to:
\begin{enumerate}
\item The expenditure function.
\item The profit function.
\item The production function.
\item The marginal utility function.
\end{enumerate}
\newpage
\item If the cost function is $C(w,r,Q)$, Shephard's Lemma states:
\begin{enumerate}
\item $\displaystyle \frac{\partial C}{\partial Q} = MC$
\item $\displaystyle \frac{\partial C}{\partial w} = L^*$ and $\frac{\partial C}{\partial r} = K^*$
\item $\displaystyle \frac{\partial C}{\partial w} = -L^*$
\item $\displaystyle \frac{\partial C}{\partial r} = -K^*$
\end{enumerate}

\item The integral $\displaystyle \int_0^T I(t)\, dt$ represents:
\begin{enumerate}
\item Total interest earned.
\item Total depreciation of capital.
\item Total accumulated investment over $[0,T]$.
\item Marginal change in investment.
\end{enumerate}

\item A first-order linear differential equation takes the form:
\begin{enumerate}
\item $\displaystyle \frac{dx}{dt} = a x^2 + b$
\item $\displaystyle \frac{dx}{dt} + P(t)x = Q(t)$
\item $\displaystyle \frac{d^2x}{dt^2} + ax = 0$
\item $\displaystyle x = \int ax \, dt$
\end{enumerate}

\item If capital evolves as $\frac{dK}{dt}=I-\delta K$, the long-run steady-state level of capital is:
\begin{enumerate}
\item $K^* = I$
\item $K^* = \delta I$
\item $K^* = \frac{1}{I}$
\item $K^* = \frac{I}{\delta}$
\end{enumerate}

\end{enumerate}

---

\section*{Part B — Short Answer Problems}

\textbf{Question 1 — Envelope Theorem (4 marks)}\\
Consider the constrained optimization problem:
\[
\max_{x,y} \; U(x,y,a) \quad \text{s.t. } g(x,y,a)=0,
\]
with Lagrangian:
\[
\mathcal{L}=U(x,y,a)+\lambda[g(x,y,a)].
\]

\begin{enumerate}
\item Derive the expression for the value function derivative $\frac{dV}{da}$ using the Envelope Theorem.  
\item Explain (in one short paragraph) why the indirect effect through $x^*$ and $y^*$ drops out.  
\item Provide an economic interpretation of the term $-\lambda \frac{\partial g}{\partial a}$.  
\end{enumerate}

\vspace{1em}

\textbf{Question 2 — Duality and Shephard’s Lemma (4 marks)}\\
Suppose the cost function for a firm is:
\[
C(w,r,Q) = wL^*(Q) + rK^*(Q),
\]
where $w$ is the wage and $r$ is the rental rate of capital.

\begin{enumerate}
\item Use Shephard’s Lemma to derive the firm's conditional factor demands $L^*(Q)$ and $K^*(Q)$.  
\item Briefly explain why cost-minimizing behavior implies a dual relationship between cost and production.  
\item Show that marginal cost equals $\displaystyle \frac{\partial C}{\partial Q}$.  
\end{enumerate}

\vspace{1em}

\textbf{Question 3 — Integration and Stock–Flow Dynamics (4 marks)}\\
Let capital evolve according to the equation:
\[
\frac{dK}{dt}=I(t)-\delta K(t).
\]

\begin{enumerate}
\item Write the integrating factor and show all steps needed to solve for $K(t)$.  
\item Derive the long-run (steady-state) value of $K$.  
\item Provide an economic interpretation of the terms:
\[
K(0)e^{-\delta t}, \qquad \frac{1}{\delta}\int_0^t e^{-\delta (t-s)}I(s)ds.
\]
\end{enumerate}

\vspace{1em}

\textbf{Question 4 — Present Value Calculation (4 marks)}\\
A consumer receives income according to:
\[
Y(t)=Y_0 e^{gt},
\]
where $Y_0>0$ and $g$ is the income growth rate. The discount rate is $r>g$.

\begin{enumerate}
\item Compute the present value:
\[
PV=\int_0^\infty Y(t) e^{-rt} dt.
\]
Show all integration steps.  
\item State the condition needed for the present value to be finite and explain why.  
\end{enumerate}

---

\textbf{End of Quiz}

\end{document}
